\chapter{1998年上海卷（理）}
\section{填空题}

\begin{problem}
\(\lg20 + \log_{100}25 = \)\fillinblank{}.

\end{problem}

\begin{answer}
\(2\)
\end{answer}

\begin{problem}
若函数 \(y = 2\sin x + \sqrt{a} \cos x + 4 \) 的最小值为 \(1\)，则 \(a=\)\fillinblank{}.

\end{problem}

\begin{answer}
\(5\)
\end{answer}

\begin{problem}
若 \(\lim\limits_{x\to1}\frac{x^{2}+ax+3}{x^{3}+3}=2 \)，则 \(a=\)\fillinblank{}.

\end{problem}

\begin{answer}
\(4\)
\end{answer}

\begin{problem}
函数 \(f(x)=(x-1)^{\frac{1}{3}}+2 \) 的反函数是 \(f^{-1}(x)= \)\fillinblank{}.

\end{problem}

\begin{answer}
\((x-2)^3+1\)，\(\ x\in\R\)
\end{answer}

\begin{problem}
棱长为 \(2\) 的正四面体的体积为\fillinblank{}.

\end{problem}

\begin{answer}
\(\frac{2\sqrt{2}}{3}\)
\end{answer}

\begin{problem}
以直角坐标系的原点 \(O\) 为极点，\(x\) 轴的正半轴为极轴建立极坐标系，若椭圆两焦点的极坐标分别是 \(\left(1, \frac{\pi}{2}\right) \)， \(\left(1, \frac{3\pi}{2}\right) \)，长轴长是 \(4\)，则此椭圆的直角坐标方程是\fillinblank{}.

\end{problem}

\begin{answer}
\(\frac{x^2}{3}+\frac{y^2}{4}=1\)
\end{answer}

\begin{problem}
袋内装有大小相同的 \(4\) 个白球和 \(3\) 个黑球，从中任意摸出 \(3\) 个球，其中只有一个黑球的概率是\fillinblank{}.

\end{problem}

\begin{answer}
\(\frac{18}{35}\)
\end{answer}

\begin{problem}
函数 \(y=\left\{\begin{smallmatrix}2x+3,&x\leqslant0\\ x+3,&0<x\leqslant1\\ -x+5,&x>1\end{smallmatrix}\right.\) 的最大值是\fillinblank{}.

\end{problem}

\begin{answer}
\(4\)
\end{answer}

\begin{problem}
设 \(n\) 是一个自然数，\(\left(1+\frac{x}{n}\right)^{n} \) 的展开式中 \(x^{3} \) 的系数为 \(\frac{1}{16} \)，则 \(n=\)\fillinblank{}.

\end{problem}

\begin{answer}
\(4\)
\end{answer}

\begin{problem}
在数列 \(\{a_{n}\} \) 和 \(\{b_{n}\} \) 中， \(a_{1}=2 \)，且对任意自然数 \(n\)， \(3a_{n+1}-a_{n}=0 \)， \(b_{n} \) 是 \(a_{n} \) 与 \(a_{n+1} \) 的等差中项，则 \(\{b_{n}\} \) 的各项和是\fillinblank{}.

\end{problem}

\begin{answer}
\(2\)
\end{answer}

\begin{problem}
函数 \(f(x)=a^{x}\)（\(a>0\)，\(a\neq1\)）在 \([1,2] \) 中的最大值比最小值大 \(\frac{a}{2} \)，则 \(a\) 的值为\fillinblank{}.

\end{problem}

\begin{answer}
\(\frac{1}{2}\) 或 \(\frac{3}{2}\)
\end{answer}

\section{单选题}

\begin{problem}
下列函数中，周期为 \(\frac{\pi}{2}\) 的偶函数是（\quad）

\choices
    {\(y = \sin 4x \)}
    {\(y = \cos^{2}2x - \sin^{2}2x \)}
    {\(y = \tan 2x \)}
    {\(y = \cos 2x \)}

\end{problem}

\begin{answer}
B
\end{answer}

\begin{problem}
若 \(0 < a < 1 \)，则函数 \(y = \log_{a}(x + 5) \) 的图象不经过（\quad）

\choices
    {第一象限}
    {第二象限}
    {第三象限}
    {第四象限}

\end{problem}

\begin{answer}
A
\end{answer}

\begin{problem}
在下列命题中，假命题是（\quad）

\choices
    {若平面 \(\alpha \) 内的一条直线 \(l\) 垂直于平面 \(\beta \) 内的任一直线，则 \(\alpha \perp \beta \)}
    {若平面 \(\alpha \) 内的任一直线平行于平面 \(\beta \)，则 \(\alpha \parallel \beta \)}
    {若平面 \(\alpha \perp \beta \)，任取直线 \(l \subset \alpha \)，则必有 \(l \perp \beta \)}
    {若平面 \(\alpha \parallel \beta\)，任取直线 \(l \subset \alpha\)，则必有 \(l \parallel \beta\)}

\end{problem}

\begin{answer}
C
\end{answer}

\begin{problem}
设全集为 \(\R\)，\(A = \{x \mid x^2 - 5x - 6 > 0\}\)，\(B = \{x \mid |x - 5| < a\}\)（\(a\) 是常数），且 \(11 \in B\)，则（\quad）

\choices
    {\(\overline{A} \cup B = \R\)}
    {\(A \cup \overline{B} = \R\)}
    {\(\overline{A\cup B} = \R\)}
    {\(A \cup B = \R\)}

\end{problem}

\begin{answer}
D
\end{answer}

\begin{problem}
设 \(a\)，\(b\)，\(c\) 分别是 \(\triangle ABC\) 中 \(\angle A\)，\(\angle B\)，\(\angle C\) 所对边的边长，则直线 \(\sin A \cdot x + ay + c = 0\) 与 \(bx - \sin B \cdot y + \sin C = 0\) 的位置关系是（\quad）

\choices
    {平行}
    {重合}
    {垂直}
    {相交但不垂直}

\end{problem}

\begin{answer}
C
\end{answer}

\section{解答题}

\begin{problem}
设 \(\alpha\) 是第二象限角，\(\sin \alpha = \frac{3}{5}\)，求 \(\sin \left(\frac{37}{6} \pi - 2\alpha\right)\) 的值.

\end{problem}

\begin{solution}
因为 \(\alpha\) 是第二象限角，\(\sin\alpha=\frac{3}{5}\)，所以 \(\cos\alpha=-\frac{4}{5}\).

于是 \(\sin2\alpha=2\sin\alpha\cos\alpha=-\frac{24}{25}\)，\(\cos2\alpha=1-2\sin^2\alpha=\frac{7}{25}\).
所以
\[
\sin\left(\frac{37}{6}\pi-2\alpha\right)
=\sin\left(\frac{\pi}{6}-2\alpha\right)
=\sin\frac{\pi}{6}\cos2\alpha-\cos\frac{\pi}{6}\sin2\alpha
=\frac{7+24\sqrt{3}}{50}
\]
\end{solution}

\begin{problem}
已知向量 \(\overrightarrow{OZ}\) 所表示的复数 \(z\) 满足 \((z - 2)\i = 1 + \i\)，将 \(\overrightarrow{OZ}\) 绕原点 \(O\) 按顺时针方向旋转 \(\frac{\pi}{4}\) 得 \(\overrightarrow{OZ'}\)，设 \(\overrightarrow{OZ'}\) 所表示的复数为 \(z'\)，求复数 \(z' + \sqrt{2}\i\) 的辐角主值.

\end{problem}

\begin{solution}
由 \((z-2)\i=1+\i\) 得 \(z=\frac{1+\i}{\i}+2=3-\i\).
于是
\[
z'=z\left[\cos\left(-\frac{\pi}{4}\right)+\i\sin\left(-\frac{\pi}{4}\right)\right]
=(3-\i)\left(\frac{\sqrt{2}}{2}-\frac{\sqrt{2}}{2}\i\right)
=\sqrt{2}-2\sqrt{2}\i
\]
所以
\[
z'+\sqrt{2}\i=\sqrt{2}-\sqrt{2}\i
=2\left(\frac{\sqrt{2}}{2}-\frac{\sqrt{2}}{2}\i\right)
\]
设 \(z'+\sqrt{2}\i\) 的辐角主值为 \(\theta\)，则 \(\cos\theta=\frac{\sqrt{2}}{2}\)，\(\sin\theta=-\frac{\sqrt{2}}{2}\)，所以 \(\theta=\frac{7\pi}{4}\). 因而 \(\arg(z'+\sqrt{2}\i)=\frac{7\pi}{4}\).
\end{solution}

\begin{problem}
直四棱柱 \(ABCD - A_{1}B_{1}C_{1}D_{1} \) 的高为 \(6\)，底面是边长为 \(4\)， \(\angle DAB = 60^{\circ} \) 的菱形，\(AC\) 与 \(BD\) 相交于 \(O\)， \(A_{1}C_{1}\) 与 \(B_{1}D_{1}\) 相交于 \(O_{1}\)，\(E\) 是 \(O_{1}A\) 的中点.

\begin{enumerate}
\item 求二面角 \(O_{1}-BC-D \) 的大小（用反三角函数表示）.
\item 分别以射线 \(OA\)，\(OB\)，\(OO_{1}\) 为 \(x\) 轴，\(y\) 轴，\(z\) 轴的正半轴建立空间直角坐标系，求点 \(B_{1}\)，\(D_{1}\)，\(E\) 的坐标，并求异面直线 \(OB_{1}\) 与 \(D_{1}E\) 所成角的大小（用反三角函数表示）.
\end{enumerate}

\end{problem}

\begin{solution}
（1）因为 \(DD_{1}\perp\) 底面 \(ABCD\)，\(O_{1}O\parallel D_{1}D\)，所以 \(O_{1}O\perp\) 底面 \(ABCD\).

过 \(O\) 作 \(OF\perp BC\)，垂足为 \(F\)，连接 \(O_{1}F\)，根据三垂线定理得 \(O_{1}F\perp BC\)，所以 \(\angle O_{1}FO\) 是二面角 \(O_{1}-BC-D\) 的平面角.

因为 \(OC=OA=4\cos30^\circ=2\sqrt{3}\)，\(OB=OD=2\)，所以
\[
OF=\frac{OC\cdot OB}{BC}=\sqrt{3}
\]
于是
\[
\tan\angle O_{1}FO=\frac{6}{\sqrt{3}}=2\sqrt{3}
\]
所以二面角 \(O_{1}-BC-D\) 的大小为 \(\arctan 2\sqrt{3}\).

（2）由条件得 \(O\)，\(B_{1}\)，\(D_{1}\)，\(E\) 的坐标为 \(O(0,0,0)\)，\(B_{1}(0,2,6)\)，\(D_{1}(0,-2,6)\)，\(E(\sqrt{3},0,3)\).
于是 \(\overrightarrow{OB_{1}}=\{0,2,6\}\)，\(\overrightarrow{D_{1}E}=\{\sqrt{3},2,-3\}\).
记向量 \(\overrightarrow{OB_{1}}\) 与 \(\overrightarrow{D_{1}E}\) 的夹角为 \(\alpha\)，则
\[
\cos\alpha=\frac{\overrightarrow{OB_{1}}\cdot\overrightarrow{D_{1}E}}{|\overrightarrow{OB_{1}}|\cdot|\overrightarrow{D_{1}E}|}
=\frac{7\sqrt{10}}{40}
\]
由此可得异面直线 \(OB_{1}\) 与 \(D_{1}E\) 所成的角为 \(\arccos\frac{7\sqrt{10}}{40}\).
\end{solution}

\begin{problem}
\begin{enumerate}
\item 动直线 \(y = a\) 与抛物线 \(y^{2} = \frac{1}{2}(x - 2) \) 相交于 \(A\) 点，动点 \(B\) 的坐标是 \((0, 3a) \)，求线段 \(AB\) 中点 \(M\) 的轨迹 \(C\) 的方程.
\item 过点 \(D(2,0)\) 的直线 \(l\) 交上述轨迹 \(C\) 于 \(P\)，\(Q\) 两点，\(E\) 点坐标是 \((1,0)\)，若 \(\triangle EPQ\) 的面积为 \(4\)，求直线 \(l\) 的倾斜角 \(\alpha \) 的值.
\end{enumerate}

\end{problem}

\begin{solution}
（1）设 \(M\) 的坐标为 \((x,y)\)，由 \(A\) 的坐标为 \((2a^{2}+2,a)\)，\(B\) 的坐标为 \((0,3a)\) 得 \(x=a^{2}+1\)，\(y=2a\)，所以轨迹 \(C\) 的方程为 \(x=\frac{y^{2}}{4}+1\)，即 \(y^{2}=4(x-1)\).

（2）【法一】设直线 \(l\) 的方程为 \(y=k(x-2)\)，因 \(l\) 与抛物线有两个交点，故 \(k\neq0\).

代入 \(y^{2}=4(x-1)\)，得
\[
y^{2}-\frac{4}{k}y-4=0
\]
且 \(\Delta=\frac{16}{k^{2}}+16>0\) 恒成立.记这个方程的两实根为 \(y_{1}\)，\(y_{2}\)，则
\[
|PQ|=\sqrt{1+\frac{1}{k^{2}}}|y_{1}-y_{2}|
=\sqrt{1+\frac{1}{k^{2}}}\sqrt{(y_{1}+y_{2})^{2}-4y_{1}y_{2}}
=\frac{4(k^{2}+1)}{k^{2}}
\]
又点 \(E\) 到直线 \(l\) 的距离
\[
d=\frac{|k\cdot1-0-2k|}{\sqrt{k^{2}+1}}=\frac{|k|}{\sqrt{k^{2}+1}}
\]
所以
\[
S_{\triangle EPQ}=\frac{1}{2}|PQ|\cdot d=\frac{2\sqrt{k^{2}+1}}{|k|}
\]
由 \(\frac{2\sqrt{k^{2}+1}}{|k|}=4\) 得 \(k^{2}=\frac{1}{3}\)，所以 \(k=\pm\frac{\sqrt{3}}{3}\)，故 \(\alpha=\frac{\pi}{6}\) 或 \(\frac{5\pi}{6}\).

\(2^{\odot}\)【法二】设直线 \(l\) 的方程为 \(y=k(x-2)\)，代入 \(y^{2}=4(x-1)\)，得
\[
k^{2}x^{2}-(4k^{2}+4)x+4k^{2}+4=0
\]
因 \(l\) 与抛物线有两个交点，故 \(k\neq0\)，而 \(\Delta=16(k^{2}+1)>0\) 恒成立.记这个方程的两个实根为 \(x_{1}\)，\(x_{2}\)，因抛物线 \(y^{2}=4(x-1)\) 的焦点是 \(D(2,0)\)，准线是 \(x=0\)，所以
\[
|PQ|=x_{1}+x_{2}=\frac{4(k^{2}+1)}{k^{2}}
\]
其余同法一.
\end{solution}

\begin{problem}
设某物体一天中的温度 \(T\) 是时间 \(t\) 的函数： \(T(t) = at^{3} + bt^{2} + ct + d\)（\(a \neq 0\)），其中温度的单位是 \(^\circ\mathrm{C}\)，时间的单位是小时，\(t = 0\) 表示 12:00，\(t\) 取正值表示 12:00 以后. 若测得该物体在 8:00 的温度为 \(8^\circ\mathrm{C}\)，12:00 的温度为 \(60^\circ\mathrm{C}\)，13:00 的温度为 \(58^\circ\mathrm{C}\)，且已知该物体的温度在 8:00 和 16:00 有相同的变化率.

\begin{enumerate}
\item 写出该物体的温度 \(T\) 关于时间 \(t\) 的函数关系式.
\item 该物体在 10:00 到 14:00 这段时间中（包括 10:00 和 14:00），何时温度最高，并求出最高温度.
\item 如果规定一个函数 \(f(x)\) 在 \([x_1, x_2]\)（\(x_1 < x_2\)）上函数值的平均为 \(\frac{1}{x_2 - x_1} \int_{x_1}^{x_2} f(x) dx\)，求该物体在 8:00 到 16:00 这段时间内的平均温度.
\end{enumerate}

\end{problem}

\begin{solution}
（1）因为 \(T(0)=60\)，\(T(-4)=8\)，\(T(1)=58\)，\(T'(-4)=T'(4)\)，所以 \(d=60\)，\(b=0\)，\(a=1\)，\(c=-3\).

因此，温度函数 \(T(t)=t^{3}-3t+60\).

（2）\(T'(t)=3t^{2}-3=3(t-1)(t+1)\).

当 \(t\in(-2,-1)\cup(1,2)\) 时，\(T'(t)>0\)；当 \(t\in(-1,1)\) 时，\(T'(t)<0\).

所以，函数 \(T(t)\) 在 \((-2,-1)\) 上单调递增，在 \((-1,1)\) 上单调递减，在 \((1,2)\) 上单调递增，即 \(t=-1\) 是极大值点.

由于 \(T(-1)=T(2)=62\)，所以在 10:00 到 14:00 这段时间中，该物体在 11:00 和 14:00 的温度最高，最高温度为 \(62^\circ\mathrm{C}\).

（3）根据定义，平均温度为
\[
\frac{1}{4-(-4)}\int_{-4}^{4}T(t)\,\mathrm{d}t
=\frac{1}{8}\int_{-4}^{4}(t^{3}-3t+60)\,\mathrm{d}t
=60
\]
即该物体在 8:00 到 16:00 这段时间内的平均温度为 \(60^\circ\mathrm{C}\).
\end{solution}

\begin{problem}
若 \(A_n\) 和 \(B_n\) 分别表示数列 \(\{a_n\}\) 和 \(\{b_n\}\) 前 \(n\) 项的和，对任意正整数 \(n\)，\(a_n = -\frac{2n + 3}{2}\)，\(4B_n - 12A_n = 13n\).

\begin{enumerate}
\item 求数列 \(\{b_n\}\) 的通项公式.
\item 设有抛物线列 \(C_1\)，\(C_2\)，\(\cdots\)，\(C_n\)，\(\cdots\)，抛物线 \(C_n\)（\(n \in \mathbf{N}^{*}\)）的对称轴平行于 \(y\) 轴，顶点为 \((a_n, b_n)\)，且通过点 \(D_n(0, n^{2} + 1)\)，过点 \(D_n\) 且与抛物线 \(C_n\) 相切的直线斜率为 \(k_n\)，求极限 \(\displaystyle \lim\limits_{n \to \infty} \frac{k_1 + k_2 + k_3 + \cdots + k_n}{a_n b_n}\).
\item 设集合 \(X = \{x \mid x = 2a_n, n \in \mathbf{N}^{*}\}\)，\(Y = \{y \mid y = 4b_n, n \in \mathbf{N}^{*}\}\). 若等差数列 \(\{c_n\}\) 的任一项 \(c_n \in X \cap Y\)，\(c_1\) 是 \(X \cap Y\) 中的最大数，且 \(-265 < c_{10} < -125\)，求 \(\{c_n\}\) 的通项公式.
\end{enumerate}
\end{problem}

\begin{solution}
（1）当正整数 \(n\geqslant2\) 时，由 \(4B_n-12A_n=13n\) 和 \(4B_{n-1}-12A_{n-1}=13(n-1)\) 相减得 \(4b_n-12a_n=13\). 由 \(a_n=-\frac{2n+3}{2}\) 得
\(4b_n=12a_n+13=-12n-5\)，所以 \(b_n=-\frac{12n+5}{4}\).
又 \(4b_1-12a_1=13\)，可得 \(b_1=-\frac{17}{4}\). 因此 \(\{b_n\}\) 的通项公式是 \(b_n=-\frac{12n+5}{4}\).

（2）设抛物线 \(C_n\) 的方程为
\[
y=a\left(x+\frac{2n+3}{2}\right)^2-\frac{12n+5}{4}
\]
因为 \(D_n(0,n^{2}+1)\) 在此抛物线上，即得 \(a=1\)，因此，\(C_n\) 的方程为
\[
y=\left(x+\frac{2n+3}{2}\right)^2-\frac{12n+5}{4}
\]
即
\[
y=x^2+(2n+3)x+n^2+1
\]
所以 \(y'=2x+(2n+3)\)，
于是 \(D_n\) 处切线斜率 \(k_n=2n+3\).
因此
\[
\lim\limits_{n\to\infty}\frac{k_1+k_2+k_3+\cdots+k_n}{a_nb_n}
=\lim\limits_{n\to\infty}\frac{n^2+4n}{\left(-n-\frac{3}{2}\right)\left(-3n-\frac{5}{4}\right)}
=\frac{1}{3}
\]

（3）对任意正整数 \(n\)，\(2a_n=-2n-3\)，\(4b_n=-12n-5=-2(6n+1)-3\in X\)，所以 \(Y\subseteq X\)，故可得 \(X\cap Y=Y\).

\(c_1\) 是 \(X\cap Y\) 中的最大数，所以 \(c_1=-17\). 设等差数列 \(\{c_n\}\) 的公差为 \(d\)，则 \(c_{10}=-17+9d\).
由 \(-265<c_{10}<-125\)，得 \(-27\frac{5}{9}<d<-12\).
而 \(\{4b_n\}\) 是一个以 \(-12\) 为公差的等差数列，所以 \(d=-12m\)（\(m\in\mathbf{N}\)），故 \(d=-24\). 因此 \(c_n=7-24n\)（\(n\in\mathbf{N}\)）.
\end{solution}
